\section{Failure of a sum-of-squares representation for Toeplitz BW Form}\label{sec:non-sos}

In this section, we prove the following theorem.
\begin{tcolorbox}[colback=white,colframe=black!60,boxrule=0.5pt,
 arc=0pt,left=7pt,right=7pt,top=3pt,bottom=3pt]
\begin{theorem}\label{thm:negative}
There exists an integer $N_0$ such that, for every integer $N\ge N_0$,
the polynomial $F_N$ in \eqref{eq:FN} is nonnegative but is not a finite
sum of squares of real homogeneous quadratic polynomials.
An explicit sufficient value of $N_0$ is given in
\Cref{app:explicit-bound}.
\end{theorem}
\end{tcolorbox}

This section is organized as follows.  In \Cref{subsec:wedge-coordinates},
following \cite[Section~2]{laszlo2012}, we introduce the wedge coordinates
and the associated wedge Gram representation.  These coordinates will be
used throughout the paper.  In \Cref{subsec:corner}, we introduce a class
of corner-supported Toeplitz matrices.  The B\"ottcher--Wenzel forms
generated by this class will provide the non-SoS examples at sufficiently
large matrix size.

Our proof proceeds in three steps.  We first define the corresponding
corner-supported B\"ottcher--Wenzel form, denoted by $\mathcal F_m$, and
construct a convenient Gram representation for it (\Cref{subsec:corner-gram}).  Since Gram representations are not unique, we
then use the symmetries of $\mathcal F_m$ to reduce the full Gram family
to a general normal form (\Cref{lem:normal-form}).  Finally, we construct a finite matrix test which, for
sufficiently large $m$, rules out positive semidefiniteness for every
possible Gram representation of $\mathcal F_m$ (\Cref{subsec:finite-test}).  This proves that
$\mathcal F_m$ is not SoS for sufficiently large $m$, and hence, that $F_N$ is not SoS for all sufficiently large $N$.

\subsection{Wedge-coordinate and wedge Gram representation}\label{app:wedge-gram}\label{subsec:wedge-coordinates}

Following \cite[Section~2]{laszlo2012}, we use the wedge-coordinate representation of \eqref{eq:FN} throughout the paper. The BW form of Toeplitz matrices $X,Y$ is a nonnegative homogeneous quartic polynomial in the $4N-2$ real variables
\begin{equation}
    \label{origin parameter}
(x_{-(N-1)},\ldots,x_{N-1},y_{-(N-1)},\ldots,y_{N-1}) \in\R^{4N-2} 
\end{equation}
with $x,y\in\R^{2N-1}$. 
We define the wedge coordinates of $(x,y)$ by
\begin{equation}
z_{ab}=x_ay_b-x_by_a,\qquad -(N-1)\le a<b\le N-1,
    \label{wedge coordinate}
\end{equation}
where $x_a$ denotes the entry of $x$ indexed by $a$ and $y_b$ denotes the entry of $y$ indexed by $b$. Collect these coordinates, in \begin{equation}
z:=\bigl(z_{-(N-1),-(N-2)},z_{-(N-1),-(N-3)},\ldots,z_{-(N-1),N-1},
z_{-(N-2),-(N-3)},\ldots,z_{N-2,N-1}\bigr)^\top
\in\R^{\binom{2N-1}{2}}.
\end{equation}
The vector $z$ is a quadratic lifting of the original variables.  
The wedge coordinates should be viewed as a structured quadratic lifting of the original Toeplitz parameters in \eqref{origin parameter} rather than as a linear change of variables.  Each Toeplitz matrix is determined by $2N-1$ real parameters, so the pair $(X,Y)$ is described by $4N-2$ real variables.  The wedge vector
$
z=(z_{ab})_{a<b}, 
z_{ab}=x_ay_b-x_by_a,
$
has
$
\binom{2N-1}{2}
$
coordinates, each of which is quadratic in the original variables.  Thus the map
$
(x,y)\mapsto z=x\wedge y
$
lifts the quartic form $F_N(x,y)$ to a quadratic form in the wedge coordinates,
$
F_N(x,y)=z^\top Qz.
$
The wedge representation is substantially smaller than the full degree-two monomial representation in the $4N-2$ original variables, whose dimension is
$
\binom{4N-1}{2}.
$
The wedge coordinates are not algebraically independent; their quadratic relations are precisely the Pl\"ucker relations described in \Cref{lem:frame}.  Thus the wedge coordinates form a structured subcollection of the usual degree-two monomial representation rather than an invertible change of coordinates.



We call the real symmetric matrix $Q\in\mathbb S^{\binom{2N-1}{2}}$ a wedge Gram matrix of $F_N$ if
$$
        F_N(x,y)=z^\top Qz.
$$
% The polynomial variables and Gram matrices remain real.\footnote{Complex test vectors are used later: a real symmetric positive semidefinite $Q$ also satisfies $v^*Qv\ge0$ for $v\in\C^d$. This does not change the real polynomial identity $F_N=z^\top Qz$ into a sesquilinear identity in complex symbol variables.}
Such a Gram matrix is generally not unique; \Cref{lem:frame} characterizes this full Gram family.
It is well known that a polynomial is SoS if and only if it admits a positive semidefinite Gram representation \cite{parrilo2000structured,lasserre2001global}. For the B\"ottcher--Wenzel form, \Cref{lem:frame} also shows that every quadratic factor in an SoS representation lies in the wedge span. Consequently,
$$
        F_N(x,y)\ \text{is SoS}
        \quad\Longleftrightarrow\quad
        F_N(x,y)=z^\top Qz
        \ \text{for some }Q\succeq0.
$$
% More precisely, if $\tilde{v}_{xy}$ denotes the vector of bilinear monomials $x_a y_b$, then there is a fixed signed matrix $S_N$ such that $z=S_N\tilde{v}_{xy}$. Let
% $\tilde{v}_{xy}$ denotes the vector of bilinear monomials $x_a y_b$ such that 

% Note that \Cref{lem:frame} gives a general formulation
% of this reduction and characterizes all admissible Gram
% representations. The lemma shows that these are not merely some Gram
% redundancies: they describe \emph{all} Gram redundancies relevant to the
% quartic forms considered below.


\begin{lemma}[Complete characterization of wedge Gram representations]\label{lem:frame}
Let $I$ be a finite ordered set. Let \(F(x,y)\) be a real homogeneous quartic satisfying
\begin{equation}
    \label{vanish}
    F(x,0)=F(0,y)=F(x,x)=0
    \qquad
    \text{for all }x,y\in\R^I.\end{equation}
 If $F$ is a sum of squares, every quadratic factor is
a linear combination of the $z_{ab}$.
Moreover, the real symmetric corrections $\mathcal E$ satisfying
$z^\top\mathcal E z\equiv0$ are spanned by the matrices
$P_{abcd}$ specified by\footnote{Precisely, its entries at $(ab,cd),(ac,bd),(ad,bc)$ are $1,-1,1$, with the same entries at transposed positions and zero elsewhere.}
\begin{equation}\label{eq:plucker}
 z^\top P_{abcd}z
 =2(z_{ab}z_{cd}-z_{ac}z_{bd}+z_{ad}z_{bc}),
 \qquad a<b<c<d.
\end{equation}
Consequently, if $F=z^\top Gz$, then
\begin{equation}\label{eq:frame}
 F\text{ is SoS}
 \quad\Longleftrightarrow\quad
 G+\mathcal E\succeq0\text{ for some }\mathcal E\in\Span\{P_{abcd}\}.
\end{equation}
\end{lemma}
\begin{proof}

Write $F=\sum_jq_j^2$. Evaluating at $y=0$ shows that every
$q_j(x,0)$ vanishes identically, and evaluating at $x=0$ gives the
analogous conclusion. Thus $q_j=x^\top A_jy$ for a real matrix $A_j$.
The identity $q_j(x,x)=0$ implies $A_j+A_j^\top=0$, whence
$q_j=\sum_{a<b}(A_j)_{ab}z_{ab}$. The coefficient vectors of the
$q_j$ give a positive semidefinite Gram matrix. 
To determine the relation space, group products of wedges by their
multisets of labels. Monomials belonging to different multisets cannot
cancel. For a multiset with a repeated label, there is at most one
nonzero wedge product up to sign, and that product is not the zero
polynomial. For four distinct labels, the three possible products are
$z_{ab}z_{cd}$, $z_{ac}z_{bd}$ and $z_{ad}z_{bc}$. Any two are linearly
independent, as is seen by comparing their coefficients, whereas their
alternating combination in \eqref{eq:plucker} vanishes. Their relation
space is therefore exactly one-dimensional. This proves the asserted
description of the correction space. 
Any representing Gram matrix differs from $G$ by an element of this
space. Conversely, a positive semidefinite representative factors as
$R^\top R$, and the rows of $Rz$ provide the required quadratic factors.
This proves \eqref{eq:frame}.
\end{proof}

Note that Lemma~\ref{lem:frame} does not require Toeplitz structure. It applied to the Toeplitz Bottcher--Wenzel form since for the polynomial $F_N$, the three vanishing conditions
\eqref{vanish} hold.

\subsection{Toeplitz subclass that is non-SoS for large $N$}
\label{subsec:corner}

We begin by identifying a Toeplitz subclass on which the large-order
failure of SoS representability will be detected.  Fix a positive integer $m$ and
 $N\ge 2m$.  We call a Toeplitz matrix
\emph{corner-supported at depth $m$} if only the outermost $m$
diagonals on each side are retained and all Toeplitz coefficients in
the central band are set to zero.
The following small example illustrates the geometry.
Take $N=6$ and $m=2$, a corner-supported Toeplitz matrix with depth $2$ has the form
\begin{equation}\label{eq:corner-example-matrix}
X_{(N=6,m=2)}
=
\begin{pmatrix}
0&0&0&0&x_{-4}&x_{-5}\\
0&0&0&0&0&x_{-4}\\
0&0&0&0&0&0\\
0&0&0&0&0&0\\
x_4&0&0&0&0&0\\
x_5&x_4&0&0&0&0
\end{pmatrix}.
\end{equation}
Note that this example is only to illustrate the corner
geometry.  At this small depth $m=2$, the corresponding BW form is
still SoS; the failure of SoS representability occurs only for sufficiently large corner
depth. 

More generally, let $X$ and $Y$ be corner-supported Toeplitz matrices with at depth $m$, then we denote the indices of the retained entry (non-zero band) of Toeplitz labels as
\[
 \Lambda_{N,m}
 :=
 \{\pm(N-1), \pm (N-2), \dotsc, \pm(N-m)\}
\]
such that $ x_a=y_a=0$ for $a \notin \Lambda_{N,m}$.
The \eqref{eq:FN} with respect to in corner-supported $X,Y$ is the restriction to the the quartic
\begin{equation}\label{eq:Bm-definition}
 \mathcal F_m
 :=
 F_N(x,y)\Big|_{
 x_a=y_a=0\;(a\notin\Lambda_{N,m})}.
\end{equation}

\begin{lemma}[Properties of the corner-supported Toeplitz BW form]\label{lem:corner}
Let $X$ and $Y$ be corner-supported Toeplitz matrices of size $N$ and
depth $m$, with $N\ge 2m$, and let $\mathcal F_m$ be the restricted
B\"ottcher--Wenzel form defined in \eqref{eq:Bm-definition}.  Let
$
 z=(z_{ab})_{a<b},
 z_{ab}:=x_a y_b-x_b y_a,
$
be the full wedge vector defined in \eqref{wedge coordinate}, and let
\begin{equation}
\label{restricted wedge coordinate}
 z_r:=(z_{ab})_{
 a,b\in\Lambda_{N,m},\,a<b}
 \in\R^{\binom{2m}{2}}
\end{equation}
be its restriction to the retained labels.
Then the following hold.
\begin{enumerate}
\item
If
$
        F_N=z^\top Gz
$
is any Gram representation in the full wedge coordinates, then
\[
        \mathcal F_m=z_r^\top G_rz_r,
\]
where $G_r$ is the principal submatrix of $G$ indexed by the retained
wedge coordinates.

\item
After identifying the retained labels by their side and depth
$p=0,\ldots,m-1$, the polynomial $\mathcal F_m$ depends only on $m$ and
is independent of $N$.

\item
If $\mathcal F_m$ is not SoS, then $F_N$ is not SoS for every
$N\ge2m$.
\end{enumerate}
\end{lemma}

\begin{proof}
For the first statement, setting $x_a=y_a=0$ for $a\notin\Lambda_{N,m}$ annihilates every wedge containing a non-retained label. Hence $z^\top Gz$ restricts directly to $z_r^\top G_rz_r$.

For the second statement, let $U_a$ be the shift matrix $(U_a)_{ij}=\one_{i-j=a}$, $0\le i,j<N$. Distinct shifts are Frobenius orthogonal and $\|U_a\|_F^2=N-|a|$. The weighted Lagrange identity \cite[Lemma~10]{laszlo2012} and bilinearity of the commutator give
\begin{equation}\label{eq:literal-gram}
 F_N=
 2\sum_{a<b}(N-|a|)(N-|b|)z_{ab}^2
 -
 \Bigl\|\sum_{a<b}z_{ab}[U_a,U_b]\Bigr\|_F^2.
\end{equation}
For a retained label, $N-|\pm(N-1-a)|=a+1$. Shifts on the same side commute, while for $C_{ab}:=[U_{-(N-1-a)},U_{N-1-b}]$ one has, for $N\ge2m$,
\begin{equation}\label{eq:commutator-overlap}
 \langle C_{ab},C_{cd}\rangle_F
 =
 2\one_{a-b=c-d}\min(a+1,b+1,c+1,d+1).
\end{equation}
Thus every coefficient of the restricted quartic depends only on the depths $a,b,c,d$, and hence only on $m$.

Finally, linear restriction preserves sums of squares. Therefore an SoS representation of $F_N$ would restrict to an SoS representation of $\mathcal F_m$, proving the last statement by contraposition.
\end{proof}

\subsubsection{Gram representation of the corner-supported Toeplitz BW form}
\label{subsec:corner-gram}

In this subsection, we study the Gram representations of the
corner-supported Toeplitz B\"ottcher--Wenzel form $\mathcal B$.  We
first construct one fixed Gram representation with a particularly simple
block structure, and then derive a general normal form that captures all
positive semidefinite Gram representations under the reduced restricted
wedge coordinates.

To obtain the structured Gram representation, we first rearrange the
restricted wedge vector as
$
z_r=(u^\top,v^\top,w^\top)^\top,
$
where
$
u,v\in\R^{\binom m2}
$
and
$
w\in\R^{m^2}.
$
The entries are defined by
\[
u_{ab}:=x_ay_b-x_by_a,\quad a<b<0,\qquad
v_{ab}:=x_ay_b-x_by_a,\quad 0<a<b, \qquad
w_{ab}:=x_ay_b-x_by_a,\quad a<0<b.
\]
We refer to $u$ and $v$ as the two \emph{pure wedge} families, since the
two labels in each wedge lie on the same side of the main diagonal, and
to $w$ as the \emph{mixed wedge} family, since its two labels lie on
opposite sides. With this ordering, the restricted Gram representation takes a highly
regular block form.  The structure is already visible in the case
$m=2$.



\begin{example}[Gram representation for a corner-supported Toeplitz BW form with $m=2$]\label{ex:gram-m2}
In~\eqref{eq:corner-example-matrix}, with the ordering described above, the corresponding restricted wedge vector and baseline Gram matrix are
\[
z_r=
\left(
\begin{array}{c}
u_{-5,-4}\\
v_{4,5}\\
w_{-5,5}\\
w_{-5,4}\\
w_{-4,5}\\
w_{-4,4}
\end{array}
\right)
\qquad
\begin{array}{c}
\left.\vphantom{\begin{array}{c}u_{-5,-4}\end{array}}\right\}\ \text{$-$ pure}\\[-0.2ex]
\left.\vphantom{\begin{array}{c}v_{4,5}\end{array}}\right\}\ \text{$+$ pure}\\[-0.2ex]
\left.\vphantom{\begin{array}{c}w_{-5,5}\\ w_{-5,4}\\ w_{-4,5}\\ w_{-4,4}\end{array}}\right\}\ \text{mixed} \\[-0.2ex]
\end{array}
\qquad
G_2=
\left(
\begin{array}{c|c|cccc}
4&2&0&0&0&0\\ \hline
2&4&0&0&0&0\\ \hline
0&0&0&0&0&0\\
0&0&0&2&-2&0\\
0&0&0&-2&2&0\\
0&0&0&0&0&4
\end{array}
\right).
\]
Here the mixed coordinates are ordered by corner depth,
$
(w_{00},w_{01},w_{10},w_{11})
=
(w_{-5,5},w_{-5,4},w_{-4,5},w_{-4,4}).
$
Thus,
\[
\mathcal B_2=z_r^\top G_2z_r
=
2(u_{-5,-4}+v_{4,5})^2
+2u_{-5,-4}^2
+2v_{4,5}^2
+2(w_{-5,4}-w_{-4,5})^2
+4w_{-4,4}^2.
\]
\end{example}
Example~\ref{ex:gram-m2} shows that this ordering arranges the restricted Gram matrix into a particularly simple block structure. The next lemma proves that the same structure holds for every corner depth $m$. 
For notational simplicity, in the rest of the paper,  we index
these three families 
$
u_{ab}:=u_{-(N-1-a),-(N-1-b)},
$
$
v_{ab}:=v_{N-1-b,N-1-a},
$
and
$
w_{ab}:=w_{-(N-1-a),N-1-b}.
$
This is only a relabelling of the same restricted wedge coordinates.


\begin{lemma}[Gram representation of the corner-supported Toeplitz BW form]\label{lem:corner-gram}
Let $\calF_m$ be a corner-supported Toeplitz B\"ottcher--Wenzel form
generated by two Toeplitz matrices $X,Y$ of size $N$ and corner depth
$m$, where $N\ge2m$.  
Collect the coordinates into
$z_r=(u^\top,v^\top,w^\top)^\top$, using lexicographic
order within each depth-indexed group.
% \footnote{The example uses the original-label ordering; the depth ordering here differs by a coordinate permutation. The corresponding Gram matrices are permuted in the same way.} 
Then
\begin{equation}\label{eq:baseline-gram}
 \mathcal F_m=z_r^\top G_mz_r,\qquad
 G_m=
 \begin{pmatrix}
 D&K&0\\
 K&D&0\\
 0&0&B
 \end{pmatrix}.
\end{equation}
Here $D,K\in\mathbb S^{\binom m2}$,
$B\in\mathbb S^{m^2}$, and
$G_m\in\mathbb S^{m(2m-1)}$ where $
\mathbb S 
$
denotes the space of real symmetric matrices and the specific entries are
\begin{equation}\label{eq:baselines}
\begin{aligned}
 (D)_{ab,cd}
 &=2(a+1)(b+1)\delta_{ac}\delta_{bd}, \quad
 (K)_{ab,cd}
  =2\one_{b-a=d-c}\min(a+1,c+1),\\
 (B)_{ab,cd}
 &=2(a+1)(b+1)\delta_{ac}\delta_{bd}
 -2\one_{a+b=c+d}\min(a+1,b+1,c+1,d+1).
\end{aligned}
\end{equation}
For $D,K$, the index pairs satisfy $0\le a<b<m$ and
$0\le c<d<m$; for $B$, all four indices range independently over
$\{0,\ldots,m-1\}$.  In particular, $D,K,B$, and hence $G_m$, are fixed by \eqref{eq:baselines} and are independent of $N$.
% \footnote{The baseline $G_m$ is one fixed choice; it is not uniquely determined by the polynomial. The admissible Gram corrections remain nonunique.}
\end{lemma}

\begin{proof}
We rearrange the restricted expression \eqref{eq:literal-gram} using
the Pl\"ucker relations, as in the Gram construction of
\cite[Sections~2 and~5]{laszlo2012}.  The raw Gram couples mixed wedges
with equal differences $a-b=c-d$.  For $a<c$ and $b<d$ satisfying this
equality, the corresponding term is
$-4\min(a+1,b+1)w_{ab}w_{cd}$.  With the orientation above, the Pl\"ucker
identity becomes
\begin{equation}\label{eq:exchange}
 u_{ac}v_{bd}+w_{ab}w_{cd}-w_{ad}w_{cb}=0,
\end{equation}
so that the mixed product is replaced by
$-u_{ac}v_{bd}+w_{ad}w_{cb}$.  Summing the resulting contributions gives
\eqref{eq:baselines}.
\end{proof}

The Gram representation in Lemma~\ref{lem:corner-gram} is not unique.  By \Cref{lem:frame}, every other representation is obtained by adding Pl\"ucker corrections.  The full Gram family is large, but the symmetries of $\mathcal F_m$
allow us to average any positive semidefinite representative into a
simple normal form.  
% The realignment map $\calR$ appearing below is the index re-pairing defined in the proof.

\begin{lemma}[General Gram normal form for the corner-supported Toeplitz BW form]
\label{lem:normal-form}
Let $\calF_m$ be a corner-supported Toeplitz B\"ottcher--Wenzel form
generated by two Toeplitz matrices $X,Y$ of size $N$ and corner depth
$m$, where $N\ge2m$.  Then the following statements hold.
\begin{itemize}

\item The Gram representation of $\calF_m$ is not unique.  If $\calF_m$
admits a positive semidefinite Gram representation, then it admits a
symmetry-reduced representative of the form
\begin{equation}\label{eq:general-gram-1}
 Q_{\rm av}=G_m+
 \begin{pmatrix}
 \mathcal E_0&\mathcal E_1&0\\
 \mathcal E_1&\mathcal E_0&0\\
 0&0&\calR\mathcal E_1
 \end{pmatrix}
 =
 \begin{pmatrix}
 D+\mathcal E_0&K+\mathcal E_1&0\\
 K+\mathcal E_1&D+\mathcal E_0&0\\
 0&0&B+\calR\mathcal E_1
 \end{pmatrix}\succeq0,
\end{equation}
where $\mathcal E_0,\mathcal E_1\in\mathbb S^{\binom m2}$,
$\mathcal E_0$ is a same-side Pl\"ucker correction acting within each
pure wedge family (equivalently, a pure four-form),
$\mathcal E_1$ is a real symmetric matrix coupling the two pure wedge
families, and $\calR$ is the realignment map induced by the corresponding
index re-pairing.\footnote{The full correction matrix is an element
$\mathcal E$ of the space characterized in Lemma~\ref{lem:frame};
$\mathcal E_0$ and $\mathcal E_1$ are the components remaining after
the symmetry reduction. The realignment operator
$
\calR:\mathbb S^{\binom m2}\to\mathbb S^{m^2}
$
is a linear map that re-pairs the four indices of a pure-wedge Gram block into the corresponding mixed-wedge indices.  In particular, it maps a correction on the pure wedge space to the induced correction on the mixed wedge space. The induced mixed correction
$\calR\mathcal E_1$ belongs to $\mathbb S^{m^2}$.  See
\Cref{app:normal-form-proof} for the detailed construction.}

\item Under the orthogonal change of coordinates
$(u,v,w)\mapsto((u+v)/\sqrt2,(u-v)/\sqrt2,w)$, the Gram representation
in \eqref{eq:general-gram-1} becomes block diagonal:
\begin{equation}\label{eq:general-gram-2}
 \widehat Q_{\rm av}
 :=U^\top Q_{\rm av}U
 =
 \begin{pmatrix}
 P_+&0&0\\
 0&P_-&0\\
 0&0&M
 \end{pmatrix},
 \tag{Block-diag Gram}
\end{equation}
where $U$ is the orthogonal matrix representing the change of coordinates 
% with $r=\binom m2$,
% \[
%  U=
%  \begin{pmatrix}
%  \frac1{\sqrt2}I_r&\frac1{\sqrt2}I_r&0\\
%  \frac1{\sqrt2}I_r&-\frac1{\sqrt2}I_r&0\\
%  0&0&I_{m^2}
%  \end{pmatrix},
%  \qquad U^\top U=I,
% \]
and
\begin{equation}\label{eq:sectors}
 P_\pm=D+\mathcal E_0\pm K\pm\mathcal E_1
 \in\mathbb S^{\binom m2},
 \qquad
 M=B+\calR\mathcal E_1\in\mathbb S^{m^2}.
\end{equation}

\item The corner-supported form satisfies
\[
 \calF_m\text{ is SoS}
 \quad\Longleftrightarrow\quad
 \text{there exist }\mathcal E_0,\mathcal E_1
 \text{ such that }Q_{\rm av}\succeq0
 \quad\Longleftrightarrow\quad
 P_+\succeq0,\quad P_-\succeq0,\quad M\succeq0.
\]
\end{itemize}
\end{lemma}

\begin{proof}
The three statements follow from the complete Pl\"ucker description of
the Gram freedom, the symmetries of $\calF_m$, and the orthogonal
decomposition of the two pure wedge families.  A detailed proof is
given in \Cref{app:normal-form-proof}.
\end{proof}

\subsubsection{A finite Gram test}\label{subsec:finite-test}

We now reduce the search over the nonunique Gram family to a finite-dimensional
matrix test; we denote the matrix as $ \mathcal G_m$.  The purpose of the construction is to separate the test matrix
into a part depending only on the fixed polynomial data and a remaining part
depending on the choice of Gram representation (i.e., $  \mathcal G_m=\mathbf L_m-\mathbf R_m $).  We will show that the first
part has a negative limiting direction, whereas the second part is too small
to compensate for it. 

\begin{lemma}[Finite Gram test]\label{lem:finite-test}
Let $V\in\C^{m(2m-1)\times s}$ be an arbitrary complex matrix and let
$\mathcal G_m$ be a complex Hermitian\footnote{$\mathbb H^s:=\{A\in\mathbb C^{s\times s}:A^*=A\}$ represents the space of complex Hermitian $s\times s$ matrices.} $s\times s$ matrix defined as
\begin{equation}\label{eq:finite-compression}
 \mathcal G_m:=V^*\widehat Q_{\rm av}V\in\mathbb H^s.
\end{equation}
If $\mathcal F_m$ is SoS, then $\mathcal G_m\succeq0$ for every $V$.  In particular,
$c^*\mathcal G_mc\ge0$ for every $c\in\C^s$.
\end{lemma}

\begin{proof}
If $\mathcal F_m$ is SoS, Lemma~\ref{lem:normal-form} gives a normal-form
Gram representation with $\widehat Q_{\rm av}\succeq0$.  Hence, for every
$\hat y\in\C^s$,
\[
 \hat y^*\mathcal G_m\hat y
 =\hat y^*V^*\widehat Q_{\rm av}V\hat y
 =(V\hat y)^*\widehat Q_{\rm av}(V\hat y)\ge0.
\]
Thus $\mathcal G_m\succeq0$, and the last assertion follows by taking
$\hat y=c$.
\end{proof}

\begin{remark}[Complex test points]
\label{rem:complex-test}
The original SoS problem and all Gram representations above remain over
the real numbers.  Complex numbers are introduced only in the test
points used below, and do not change the underlying real SoS problem.

Indeed, let $Q$ be a real symmetric positive semidefinite matrix.  For
any complex vector $z=\hat u+i\hat v$, with
$\hat u,\hat v$ real, we have
\[
 z^*Qz
 =
 (\hat u^\top-i\hat v^\top)Q(\hat u+i\hat v)
 =
 \hat u^\top Q\hat u+\hat v^\top Q\hat v\ge0,
\]
where the cross terms cancel because $Q$ is symmetric.  Thus a real
positive semidefinite Gram matrix remains positive semidefinite when
tested against complex vectors.

We may therefore choose several complex test points and arrange their
corresponding feature vectors as the columns of a complex test matrix
$\mathcal V$.  The resulting matrix
$
\mathcal V^*Q\mathcal V
$
must still be positive semidefinite whenever $Q\succeq0$.  Allowing
complex test points gives us a larger and more flexible space of test
directions, without enlarging the class of admissible Gram
representations. 
In the construction in \Cref{sec choose direction}, we  will choose complex points that approach
$(i,-i)$ along prescribed directions as $m$ increases.  The associated
test matrix develops a strictly negative limiting direction $c$.  Since any
real SoS Gram representation would remain positive semidefinite under
the same complex test, this eventually yields the contradiction used to
prove that the corner-supported form is not SoS.
\end{remark}

\begin{lemma}[Finite Gram decomposition]\label{lem:finite-decomposition}
Assume $V$ contains a special block structure, namely
\begin{equation}\label{eq:special-V}
V=
\left(
\begin{array}{ccc}
V_1^*/\sqrt2&
V_1^*/\sqrt2&
V_3^*
\end{array}
\right)^*
\in\C^{m(2m-1)\times s},
\end{equation}
where $V_1 \in \C^{\binom m2 \times s}$ and
$V_3 \in \C^{m^2\times s}$.
Then, the finite Gram matrix in \eqref{eq:finite-compression} admits the decomposition
\begin{eqnarray}\label{eq:finite-decomposition}
   \mathcal G_m &=& \mathbf L_m-\mathbf R_m
   \\ &=&
\underbrace{
 V_1^*DV_1
 +
 V_3^*
 \bigl(B+\calR(D-K)\bigr)V_3
}_{:= \mathbf L_m \text{ Gram-invariant}}
-
\underbrace{
 V_3^*(\calR P_-)V_3
}_{:= \mathbf R_m \text{ Gram-dependent}} .
\end{eqnarray}
Here, \emph{Gram-invariant} means that $\mathbf L_m$ depends only on the fixed polynomial data $D,K,B$ and is independent of the particular Gram representation, whereas \emph{Gram-dependent} means that $\mathbf R_m$ depends on the chosen Gram representative through $P_-$.
\end{lemma}

\begin{proof}
From the block diagonal form of $\widehat Q_{\rm av}$,
\[
\begin{aligned}
\mathcal G_m
=
\left(
\begin{array}{ccc}
V_1^*/\sqrt2&
V_1^*/\sqrt2&
V_3^*
\end{array}
\right)
\left(
\begin{array}{ccc}
P_+&0&0\\
0&P_-&0\\
0&0&M
\end{array}
\right)
\left(
\begin{array}{c}
V_1/\sqrt2\\
V_1/\sqrt2\\
V_3
\end{array}
\right)=
V_1^*\frac{P_++P_-}{2}V_1
+
V_3^*MV_3.
\end{aligned}
\]
Using \eqref{eq:sectors},
$
(P_++P_-)/2=D+\mathcal E_0
$
and
$
M=B+\calR\mathcal E_1.
$
Moreover,
$
P_-=D-K+\mathcal E_0-\mathcal E_1
$
 and $\calR \mathcal E_0=-\mathcal E_0$\footnote{By definition, realignment exchanges the middle two indices,
$
 (\calR \mathcal E_0)_{ac,bd}
 =(\mathcal E_0)_{ab,cd}
 =\Omega_{abcd}.
$
Since $\Omega$ is totally antisymmetric,
$
 \Omega_{acbd}=-\Omega_{abcd}.
$
Equivalently,
$
 (\calR \mathcal E_0)_{ac,bd}
 =-(\mathcal E_0)_{ac,bd}.
$
Hence $\calR \mathcal E_0=-\mathcal E_0$.}.
% \footnote{This is an equality of scalar generating expressions, not of pure- and mixed-space matrices, which have different sizes. Realignment exchanges the middle two arguments of the totally alternating generating expression, changing its sign.}. 
At the level of the corresponding finite evaluations,
this gives
$
V_1^*\mathcal E_0V_1=-V_3^*(\calR \mathcal E_0)V_3.
$
Hence
\[
\begin{aligned}
\mathcal G_m
&=
V_1^*DV_1
+
V_3^*BV_3
-
V_3^*\calR(\mathcal E_0-\mathcal E_1)V_3\\
&=
V_1^*DV_1
+
V_3^*
\bigl(B+\calR(D-K)\bigr)V_3
-
V_3^*(\calR P_-)V_3,
\end{aligned}
\]
which proves \eqref{eq:finite-decomposition}.
\end{proof}

The importance of Lemma~\ref{lem:finite-decomposition} is that
$\mathbf L_m$ depends only on the fixed matrices $D,K,B$, and hence is
independent of the chosen Gram representation.  The entire remaining
dependence on the nonunique Gram representation is contained in
$\mathbf R_m$.  We will choose a finite family of points for which
$\mathbf L_m$ develops a negative direction, while the contribution of
$\mathbf R_m$ vanishes after the natural scaling.

\subsubsection{Choice of the finite test points}
\label{sec choose direction}

We now choose a specific family of test matrices $V_m$ and a direction
$c$ for which the finite compression
$
\mathcal G_m=V_m^*\widehat Q_{\rm av}V_m
$
cannot remain positive semidefinite when $m$ is sufficiently large.
% The construction is designed so that, in the decomposition
% $
% \mathcal G_m=\mathbf L_m-\mathbf R_m,
% $
% the Gram-invariant part $\mathbf L_m$ develops a strictly negative
% direction, while the Gram-dependent part $\mathbf R_m$ becomes
% asymptotically negligible in the same direction.  Since every SoS
% representation would require $\mathcal G_m\succeq0$, such a choice of
% $V_m$ and $c$ rules out all positive semidefinite Gram representations
% of $\mathcal F_m$.

Following the definition of $V$ in \eqref{eq:special-V}. We next choose $V$ so that the three diagonal blocks of
$\widehat Q_{\rm av}$ are tested on the following
coordinates.    For points
$\mathsf x_\ell=(\xi_\ell,\eta_\ell)\in\C^2$, $1\le\ell\le s$, define
\begin{equation}\label{eq:matched-features}
V_1:=
\left[
 \bigl(\xi_\ell^a\eta_\ell^b-\xi_\ell^b\eta_\ell^a\bigr)_{0\le a<b<m}
\right]_{\ell=1}^s
\in\C^{\binom m2 \times s},
\qquad
V_3:=
\left[
 \bigl(\xi_\ell^a\eta_\ell^b\bigr)_{0\le a,b<m}
\right]_{\ell=1}^s
\in\C^{m^2\times s},
\end{equation}
For each $m$, define
\begin{equation}\label{eq:test-nodes}
 \mathsf x_\ell^{(m)}
 =
 \bigl(\xi_\ell^{(m)},\eta_\ell^{(m)}\bigr)
 :=
 \left(
 ie^{-a_\ell/\sqrt m},
 -ie^{-\bar a_\ell/\sqrt m}
 \right),
 \qquad 1\le\ell\le s.
\end{equation}
where the parameters begin with
\[
\begin{alignedat}{6}
a_1&=1+i,      &\qquad a_2&=1-i,      &\qquad a_3&=2+i,      &\qquad a_4&=2-i,      &\qquad a_5&=1+2i,      &\qquad a_6&=1-2i,\\
a_7&=2+2i,     &\qquad a_8&=2-2i,     &\qquad a_9&=1+3i,     &\qquad a_{10}&=1-3i,  &\qquad a_{11}&=2+3i,  &\qquad a_{12}&=2-3i,\ \dotsc
\end{alignedat}
\]
and the pattern continues in conjugate pairs on the two vertical lines
$\Re a=1$ and $\Re a=2$. 

Note that $\mathsf x_\ell^{(m)}\to(i,-i)$ as $m\to\infty$.  From now on
$V_1,V_3,V,\mathbf L_m$ and $\mathbf R_m$ refer to the points
in \eqref{eq:test-nodes}.
The choice of $\Re a=1,2$ is not the only negative example.  Additional
families may also be included in the finite test (see
Appendix~\ref{C:rem:alternative-test-families}).

\begin{proposition}[Negative limit of the Gram-invariant part]
\label{lem:negative-limit}
Let $\mathbf L_m$ and $\mathbf R_m$ be defined as in
\Cref{lem:finite-decomposition}.  Fix the $s$ test points in
\eqref{eq:test-nodes}.  Then the following statements hold.
\begin{enumerate}

\item As $m\to\infty$,
$
\mathsf x_\ell^{(m)}\to(i,-i)
$
and
\begin{equation}\label{eq:L-limit}
        \frac{1}{m^2}\mathbf L_m\longrightarrow\mathbf L_\star,
\end{equation}
where
$
\mathbf L_\star\in\mathbb S^s
$
is a fixed real symmetric $s\times s$ matrix, independent of $m$.

\item There exist a constant $\gamma>0$ and an integer $m_0>0$ such that $
        \lambda_{\min}(\mathbf L_\star)<-\gamma<0. $
Let $c_\star\in\R^s$ be a corresponding real unit eigenvector.  Since
$\mathbf L_\star$ is independent of $m$, the same $c_\star$ is used for
every $m\ge m_0$.  

\item There exists an integer $m_1\ge m_0$ such that, for every
$m\ge m_1$, define
$
\widetilde c_m:=c_\star/m, 
$ then
\begin{equation}\label{eq:L-negative-limit}
 \widetilde c_m^*\mathbf L_m\widetilde c_m
 < -\frac{\gamma}{2}<0 .
\end{equation}

\item If $\widehat Q_{\rm av}\succeq0$, there exists an integer
$m_2\ge m_0$ such that, for every $m\ge m_2$,
\begin{equation}\label{eq:R-vanishes}
 \left|
 \widetilde c_m^*\mathbf R_m\widetilde c_m
 \right|
 <\frac{\gamma}{4},
\end{equation}
uniformly over all positive semidefinite normal-form Gram
representations.
\end{enumerate}
\end{proposition}

\begin{proof}
We prove the four statements in order. 
For the first statement, the entries of $\mathbf L_m$ are explicit finite
sums determined by $D,K,B$ and the points $\mathsf x_\ell^{(m)}$.  Since
$
\mathsf x_\ell^{(m)}\to(i,-i)
$
at rate $m^{-1/2}$, the corresponding expressions have a leading term
of order $m^2$.  For $1\le\ell,h\le s$, define
\[
 A_{\ell h}:=
 |a_\ell+\overline{a_h}|^2,
 \qquad
 C_{\ell h}:=
 4\,\Re(a_\ell)\Re(a_h).
\]
Expanding the leading term gives
\[
 \frac{1}{m^2}(\mathbf L_m)_{\ell h}
 \longrightarrow
 2\left(
 \frac{2}{A_{\ell h}^2}
 +\frac{1}{C_{\ell h}^2}
 -\frac{1}{A_{\ell h}C_{\ell h}}
 \right)
 =:(\mathbf L_\star)_{\ell h}.
\]
Thus
$
 \mathbf L_\star:=
 \left[
 2\left(
 \frac{2}{A_{\ell h}^2}
 +\frac{1}{C_{\ell h}^2}
 -\frac{1}{A_{\ell h}C_{\ell h}}
 \right)
 \right]_{\ell,h=1}^s
 \in\mathbb S^s
$
is a fixed real symmetric $s\times s$ matrix, independent of $m$.  Since
$s$ is fixed, entrywise convergence implies convergence in any matrix
norm, proving
$
        \frac{1}{m^2}\mathbf L_m\longrightarrow\mathbf L_\star.
$
% The complete calculation and finite-tail estimates are given in
% Appendix~\ref{app:quantitative}.

\noindent 
\textbf{For the second statement}, the fixed two-family choice of test points
determines $\mathbf L_\star$ completely.  Appendix~
\ref{app:negative-witness} gives an exact finite certificate showing that
$\mathbf L_\star$ has a strictly negative eigenvalue.  Hence there exists
$\gamma>0$ such that
\[
        \lambda_{\min}(\mathbf L_\star)<-\gamma<0.
\]
Let $c_\star\in\R^s$ be a corresponding real unit eigenvector:
$
 \mathbf L_\star c_\star
 =\lambda_{\min}(\mathbf L_\star)c_\star, 
 c_\star^*c_\star=1.
$
The matrix $\mathbf L_\star$, the constant $\gamma$, and the vector
$c_\star$ are fixed and independent of $m$.  Choose $m_0$ sufficiently
large that the asymptotic estimates below are valid, and define
$
\widetilde c_m:=c_\star/m
$
for every $m\ge m_0$.

\noindent
\textbf{For the third statement,}
\[
 \widetilde c_m^*\mathbf L_m\widetilde c_m
 =
 c_\star^*
 \left(\frac{\mathbf L_m}{m^2}\right)c_\star
 \longrightarrow
 c_\star^*\mathbf L_\star c_\star
 =
 \lambda_{\min}(\mathbf L_\star)
 <-\gamma.
\]
Therefore there exists an integer $m_1\ge m_0$ such that, for every
$m\ge m_1$,
$
 \widetilde c_m^*\mathbf L_m\widetilde c_m
 < -\frac{\gamma}{2}<0.
$

\noindent
\textbf{For the fourth statement}, suppose that
$\widehat Q_{\rm av}\succeq0$.  Positivity of the mixed block first
forces the low-order mixed directions into its kernel.  The resulting
improvement at the conjugate test points can then be continued to the
realigned points appearing in $\mathbf R_m$.  This yields
\[
 \sup_{\widehat Q_{\rm av}\succeq0}
 \left|
 c_\star^*\frac{\mathbf R_m}{m^2}c_\star
 \right|
 \longrightarrow0.
\]
The detailed argument is given in Appendix~\ref{app:quantitative}.
Hence there exists an integer $m_2\ge m_0$ such that, for every
$m\ge m_2$ and every positive semidefinite normal-form Gram
representation,
$
 \left|
 \widetilde c_m^*\mathbf R_m\widetilde c_m
 \right|
 =
 \left|
 c_\star^*\frac{\mathbf R_m}{m^2}c_\star
 \right|
 <\frac{\gamma}{4}.
$
This completes the proof.
\end{proof}

% \begin{proof}
% We prove the four statements in order. 
% For the first statement, the entries of $\mathbf L_m$ are explicit finite
% sums determined by $D,K,B$ and the fixed test points
% $\mathsf x_\ell^{(m)}$.  Since
% $
% \mathsf x_\ell^{(m)}\to(i,-i)
% $
% at rate $m^{-1/2}$, put, for $1\le\ell,h\le s$,
% \[
%  A_{\ell h}:=
%  |a_\ell+\overline{a_h}|^2,
%  \qquad
%  C_{\ell h}:=
%  4\,\Re(a_\ell)\Re(a_h).
% \]
% Using
% $
% 1-e^{-z/\sqrt m}=z/\sqrt m+O(m^{-1}),
% $
% the leading terms of the sums defining $\mathbf L_m$ give
% \[
%  \frac{1}{m^2}(\mathbf L_m)_{\ell h}
%  \longrightarrow
%  2\left(
%  \frac{2}{A_{\ell h}^2}
%  +\frac{1}{C_{\ell h}^2}
%  -\frac{1}{A_{\ell h}C_{\ell h}}
%  \right)
%  =:(\mathbf L_\star)_{\ell h}.
% \]
% Since $s$ is fixed,
% entrywise convergence implies convergence in any matrix norm.  Hence
% \[
%         \frac{1}{m^2}\mathbf L_m\longrightarrow\mathbf L_\star,
% \]
% where $\mathbf L_\star\in\mathbb S^s$ is a fixed real symmetric matrix,
% independent of $m$.
% % \footnote{More details of the the tail estimate is given in \Cref{B:lem:tail}}.

% \noindent \textbf{For the second statement,} $\mathbf L_\star$ is completely determined by
% the fixed two-family set of test speeds.  Appendix~
% \ref{app:negative-witness} gives an exact rational certificate showing
% that
% \[
%         \lambda_{\min}(\mathbf L_\star)<-2^{-18}<0.
% \]
% Thus, for example, we may take
% $
% \gamma=2^{-18}.
% $
% Let $c_\star\in\R^s$ be a corresponding real unit eigenvector,
% \[
%  \mathbf L_\star c_\star
%  =\lambda_{\min}(\mathbf L_\star)c_\star,
%  \qquad
%  c_\star^*c_\star=1.
% \]
% The matrix $\mathbf L_\star$, the constant $\gamma$, and the vector
% $c_\star$ are fixed and independent of $m$.  Choose an integer $m_0$
% large enough that the asymptotic estimates below apply, and define
% $
% \widetilde c_m:=c_\star/m
% $
% for $m\ge m_0$.

% \noindent \textbf{For the third statement,}
% \[
%  \widetilde c_m^*\mathbf L_m\widetilde c_m
%  =
%  c_\star^*
%  \left(\frac{\mathbf L_m}{m^2}\right)c_\star
%  \longrightarrow
%  c_\star^*\mathbf L_\star c_\star
%  =
%  \lambda_{\min}(\mathbf L_\star)
%  <-\gamma.
% \]
% Hence there exists an integer $m_1\ge m_0$ such that, for every
% $m\ge m_1$,
% \[
%  \widetilde c_m^*\mathbf L_m\widetilde c_m
%  < -\frac{\gamma}{2}<0.
% \]

% \noindent \textbf{For the fourth statement}, suppose that
% $\widehat Q_{\rm av}\succeq0$.  The low-total mixed directions are then
% forced into the kernel of every feasible mixed block.  This gives an
% improved estimate at the conjugate test points.  The qualitative argument
% in \Cref{B:prop:R} gives
% \[
%  \sup_{\widehat Q_{\rm av}\succeq0}
%  \left|
%  c_\star^*\frac{\mathbf R_m}{m^2}c_\star
%  \right|
%  \longrightarrow0.
% \]
% Therefore there exists an integer $m_2\ge m_0$ such that, for every
% $m\ge m_2$ and every positive semidefinite normal-form Gram
% representation,
% $
%  \left|
%  \widetilde c_m^*\mathbf R_m\widetilde c_m
%  \right|
%  =
%  \left|
%  c_\star^*\frac{\mathbf R_m}{m^2}c_\star
%  \right|
%  <\frac{\gamma}{4}.
% $
% This completes the proof.
% \end{proof}

Finally, we piece together the preceding lemmas to conclude the section and prove Theorem~\ref{thm:negative}.

\begin{proof}[Proof of Theorem~\ref{thm:negative}]
Let
$
M_0:=\max\{m_1,m_2\}
$
and suppose, for contradiction, that $\mathcal F_m$ is SoS for some
$m\ge M_0$.  By Lemma~\ref{lem:finite-test}, the corresponding finite
Gram test satisfies
$
\mathcal G_m\succeq0,
$
and therefore
$
        \widetilde c_m^*\mathcal G_m\widetilde c_m\ge0.
$
On the other hand, \Cref{lem:finite-decomposition,lem:negative-limit} give
\[
\begin{aligned}
\widetilde c_m^*\mathcal G_m\widetilde c_m
&\le
\widetilde c_m^*\mathbf L_m\widetilde c_m
+
\left|
\widetilde c_m^*\mathbf R_m\widetilde c_m
\right| <
-\frac{\gamma}{2}+\frac{\gamma}{4}
=
-\frac{\gamma}{4}<0,
\end{aligned}
\]
contradicting $\mathcal G_m\succeq0$. Therefore $\mathcal F_m$ is not SoS for every sufficiently large corner
depth $m$.  In particular, choose any integer $m_\star\ge M_0$.
By Lemma~\ref{lem:corner}, $\mathcal F_{m_\star}$ is a restriction of
$F_N$ for every $N\ge2m_\star$.  Hence $F_N$ is not SoS for every
$N\ge2m_\star$.  Taking
$
N_0:=2m_\star
$
proves the theorem.
\end{proof}
